\section{Context and Motivation}
\label{sec:motivation}

Active mechatronic devices such as powered prostheses and exoskeletons consist of multiple interacting subsystems.
These include mechanical structures, actuators, sensors, and embedded controllers.
Together, they form tightly coupled multi-domain systems.
Such systems are typical cyber-physical systems because computation and physical behavior are linked through sensing, actuation, and feedback \cite{lee_cyber_2008}.
A change in one subsystem can influence the behavior of the whole device.
For example, modifications in the control law, actuator design, or geometry of the mechanical structure can affect loads, stability, performance, and safety.
For this reason, system-level simulation is needed already in early design phases, before physical prototypes are available \cite{VirtualEngineeringSystems}.

For devices that interact with the human body, this challenge becomes even more demanding.
A prosthetic knee or an active exoskeleton is not merely a technical device.
Its function depends on the coupled dynamics of the device and the musculoskeletal system of the user.
In addition, local structural effects such as compliance, deformation, and contact may become relevant in critical operating phases.
Questions such as whether a controller remains stable during ground contact, how loads are transferred through compliant components, or how device forces interact with musculoskeletal system cannot be answered from isolated component analyses alone.
They require virtual testing environments that represent both the technical system and its biomechanical context in a single, time-synchronous simulation \cite{eberle_biomechanical_2025}.

In engineering practice, however, the required subsystem models are rarely created within a single tool.
Mechanical and structural behavior is often derived from geometry-based models in computer-aided design and finite element analysis environments.
Control-oriented, electrical, and other system-level subsystems are commonly formulated as one-dimensional equation-based models, for example in Modelica \cite{Modelica_Fritzson}.
Musculoskeletal dynamics, in turn, are modeled in specialized biomechanics software such as OpenSim \cite{OpenSim_Delp}.
These modeling environments serve different purposes, operate at different abstraction levels, and use different numerical representations.
Modern system development therefore requires not only reusable models within individual domains, but also a practical way to integrate heterogeneous models into a coherent system-level simulation workflow.

Co-simulation provides a natural integration strategy for such heterogeneous models.
Each subsystem retains its domain-specific solver and model representation, while a master algorithm orchestrates data exchange at discrete communication points \cite{gomes_co-simulation_2019-1}.
The Functional Mock-up Interface (FMI) standardizes the component interface for this purpose and is widely adopted in industry \cite{FMI2.0}.

However, FMI defines the \emph{interface} to individual simulation components, but it does not prescribe how multiple components are orchestrated into a consistent system simulation \cite{gomes_co-simulation_2019-1}.
Several challenges remain when heterogeneous models are coupled in practice.
These challenges include determining a correct execution order when component outputs depend algebraically on inputs \cite{broman_requirements_2015}, detecting and resolving circular dependencies between components \cite{sicklinger_interface_2014}, and handling discrete events such as mechanical contact or discrete mode changes during otherwise continuous operation \cite{broman_determinate_2013,StepRevision_Cremona,cremona_hybrid_2019}.

Moreover, different representations of the same physical subsystem involve different trade-offs between computational cost and the level of detail provided by the simulation.
A rigid-body model computes global kinematics and kinetics efficiently, but it does not explicitly resolve local deformation or stress.
An equivalent continuum-mechanical finite element model can resolve these local field quantities, but at significantly higher computational cost.
In system simulation, this suggests that detailed models should be activated selectively.
For example, the finite element representation may only be needed during contact phases in which local deformation is relevant, whereas a computationally inexpensive rigid-body model may be sufficient during free motion.
To enable such selective model use, the simulation framework must support switching between alternative models of the same subsystem at runtime, together with consistent state transfer across model boundaries \cite{Choi2014}.